% TOP_PROBLEM: {"problemId": "problem.polynomial-time-quantum-algorithm-for-the-dihedral-hidden-subgroup-problem", "canonicalTitle": "Polynomial-Time Quantum Algorithm for the Dihedral Hidden Subgroup Problem", "releaseRank": 216, "primaryDomain": "quantum_information_computation", "primaryDomainLabel": "Quantum information & computation", "displayStatus": "open", "releaseStatus": "open"}
% !TeX program = lualatex
% Definition and sources only; this file is not an LLM solution attempt.
\documentclass[11pt]{article}
\usepackage[margin=25mm]{geometry}
\usepackage{fontspec}
\setmainfont{Latin Modern Roman}
\usepackage{amsmath,amssymb,mathtools}
\usepackage{unicode-math}
\setmathfont{Latin Modern Math}
\usepackage{xurl,hyperref}
\hypersetup{colorlinks=true,urlcolor=blue}
\setcounter{secnumdepth}{0}
\setlength{\parindent}{0pt}
\setlength{\parskip}{5pt}
\setlength{\emergencystretch}{3em}
\begin{document}
\section*{\texorpdfstring{Polynomial-Time Quantum Algorithm for the Dihedral Hidden Subgroup Problem}{Polynomial-Time Quantum Algorithm for the Dihedral Hidden Subgroup Problem}}\label{216}
\subsection{Definitions and mathematical statement}
Use the convention
$D_N=\langle r,s:r^N=s^2=e,\ srs=r^{-1}\rangle$, a group of order $2N$. An oracle $f:D_N\to\mathcal Y$ hides $H$ when its fibers are the left cosets of $H$, and a quantum algorithm may query the oracle coherently. The target is a bounded-error algorithm outputting generators for $H$ with time and workspace both polynomial in $\log N$, under the standard succinct group encoding. Measuring polynomial complexity in $N$ instead would weaken the claim exponentially. A subexponential algorithm or a polynomial-query procedure whose nonquery processing is expensive does not meet the selected time requirement.
\par\smallskip\textit{Formulation and concept references:} \href{https://doi.org/10.1137/S009753970343603X}{[S1]}.

\subsection{Short English statement}
Can a quantum computer recover a hidden subgroup of a dihedral group using time and memory polynomial in the group's logarithmic size?
\subsection{Sources}
\begin{itemize}
\item[S1]G. Kuperberg, SIAM J. Comput. 35(1):170-188, 2005.
\url{https://doi.org/10.1137/S009753970343603X}.
\textit{Formulation source.}
\end{itemize}
\end{document}
